%% ma: L12-B11-0013
%% lop: 12
%% chuong: 4
%% bai: 11
%% dang: 12.11.9
%% loai: TN4
%% muc: VD
%% dap_an: "A"
%% nguon: "(NBV)-ĐỀ SỐ 17-MỖI NGÀY 1 ĐỀ THI 2025.pdf (D17, MỖI NGÀY 1 ĐỀ - Đề số 17), Phần I Câu 7"
%% kien_thuc: "Bài 11 (nguyên hàm), hàm số liên tục tại điểm nối"
%% trang_thai: cho-duyet
%% ly_do: "Dạng mới đề xuất 12.11.9 chờ giáo viên duyệt"
%% ghi_chu: "Đáp án gốc A đúng (giải độc lập: F(x)=2x^2-x+4 khi x<1, F(x)=x^2+x+3 khi x>=1, F(3)=15)."
\begin{ex}
Biết $F(x)$ là một nguyên hàm của hàm số $f(x)=\begin{cases}2x+1&\text{khi }x\ge1\\4x-1&\text{khi }x<1\end{cases}$ và $F(0)=4$. Giá trị của $F(3)$ bằng
\choice{\True $15$}{$19$}{$7$}{$11$}
\loigiai{Với $x<1$: $F(x)=2x^2-x+C_2$, vì $F(0)=4$ nên $C_2=4$. Với $x\ge1$: $F(x)=x^2+x+C_1$.

$F$ là nguyên hàm của $f$ nên $F$ liên tục tại $x=1$: $1^2+1+C_1=2\cdot1^2-1+4$, suy ra $C_1=3$.

Vậy $F(3)=9+3+3=15$. Chọn A.}
\end{ex}
